\usepackage{listings}
\usepackage[dvipsnames]{xcolor}


% Default style for all listings. Among other things, fixes 
\lstset{
	% --- Linenos ---
	numbersep=3pt,                  % how far the line-numbers are from the code
	numberstyle=\tiny\color{gray},
	numbers=left,                   % where to put the line-numbers
	stepnumber=1,                   % the step between two line-numbers. If it is 1 each line will be numbered
	% -- Basics --
	basicstyle=\ttfamily,
	sensitive=true,  % Case-sensitive keywords
	tabsize=4,
	breaklines=true,  % Break lines if too long
	%escapechar=`,
	%escapeinside={<@}{@>},
	% I don't know what keepspaces and columns do, but together, they fix the ugly default listings style.
	columns=fullflexible,
	keepspaces,
	showstringspaces=false,  % Spaces not shown as _
	basicstyle=\ttfamily\footnotesize,
}

% Usage: \begin{lstlisting}[language=notepadpp-xml] ...
\lstdefinelanguage{XML}{
	morestring=[b]",                       % "attribute strings"
	morecomment=[s]{!--}{--},            % <!-- comments -->
	morecomment=[s]{?}{?},               % <? xml instructions ?>
	morekeywords={id}, % tag names (customize)
	keywordstyle=\color{YellowGreen},      % element/tag names
	identifierstyle=\color{NavyBlue},  % attribute names
	stringstyle=\color{Fuchsia},    % attribute values
	commentstyle=\color{ForestGreen},  % XML comments
	moredelim=[s][\color{black}]{>}{<},
	showstringspaces=false,
	breaklines=true
}

% Usage: \begin{lstlisting}[language=notepadpp-xml] ...
\lstdefinelanguage{racketSQL}{
%	language=SQL,
%	morestring=[b]",                   % "attribute strings"
	morecomment=[s]{;;}{;;},           % <;; comments ;;>
	morekeywords={ABSOLUTE,ACTION,ADD,ALLOCATE,ALTER,ARE,AS,ASSERTION,%
		AT,BEGIN,BETWEEN,BIT_LENGTH,BOTH,BY,CASCADE,CASCADED,CASE,CAST,%
		CATALOG,CHAR_LENGTH,CHARACTER_LENGTH,CLUSTER,COALESCE,%
		COLLATE,COLLATION,COLUMN,CONNECT,CONNECTION,CONSTRAINT,%
		CONSTRAINTS,CONVERT,CORRESPONDING,CREATE,CROSS,CURRENT_DATE,%
		CURRENT_TIME,CURRENT_TIMESTAMP,CURRENT_USER,DAY,DEALLOCATE,%
		DEC,DEFERRABLE,DEFERED,DESCRIBE,DESCRIPTOR,DIAGNOSTICS,%
		DISCONNECT,DOMAIN,DROP,ELSE,END,EXEC,EXCEPT,EXCEPTION,EXECUTE,%
		EXTERNAL,EXTRACT,FALSE,FIRST,FOREIGN,FROM,FULL,GET,GLOBAL,%
		GRAPHIC,HAVING,HOUR,IDENTITY,IMMEDIATE,INDEX,INITIALLY,INNER,%
		INPUT,INSENSITIVE,INSERT,INTO,INTERSECT,INTERVAL,%
		ISOLATION,JOIN,KEY,LAST,LEADING,LEFT,LEVEL,LIMIT,LOCAL,LOWER,%
		MATCH,MINUTE,MONTH,NAMES,NATIONAL,NATURAL,NCHAR,NEXT,NO,NOT,NULL,%
		NULLIF,OCTET_LENGTH,ON,ONLY,ORDER,ORDERED,OUTER,OUTPUT,OVERLAPS,%
		PAD,PARTIAL,POSITION,PREPARE,PRESERVE,PRIMARY,PRIOR,READ,%
		RELATIVE,RESTRICT,REVOKE,RIGHT,ROWS,SCROLL,SECOND,SELECT,SESSION,%
		SESSION_USER,SIZE,SPACE,SQLSTATE,SUBSTRING,SYSTEM_USER,%
		TABLE,TEMPORARY,THEN,TIMEZONE_HOUR,%
		TIMEZONE_MINUTE,TRAILING,TRANSACTION,TRANSLATE,TRANSLATION,TRIM,%
		TRUE,UNIQUE,UNKNOWN,UPPER,USAGE,USING,VALUE,VALUES,%
		VARGRAPHIC,VARYING,WHEN,WHERE,WRITE,YEAR,ZONE,%
		AND,ASC,avg,CHECK,COMMIT,count,DECODE,DESC,DISTINCT,GROUP,IN,% FF
		LIKE,NUMBER,ROLLBACK,SUBSTR,sum,VARCHAR2,% FF
		MIN,MAX,UNION,UPDATE,% RF
		ALL,ANY,CUBE,CUBE,DEFAULT,DELETE,EXISTS,GRANT,OR,RECURSIVE,% DJ
		ROLE,ROLLUP,SET,SOME,TRIGGER,VIEW},% DJ,
	keywordstyle=\bfseries,      % element/tag names
%	identifierstyle=\color{NavyBlue},  % attribute names
	stringstyle=\color{Fuchsia},    % attribute values
	commentstyle=\color{ForestGreen},  % XML comments
%	moredelim=[s][\color{black}]{>}{<},
%	showstringspaces=false,
	breaklines=true
	}

% Usage: \begin{lstlisting}[language=racket] ...
\lstdefinelanguage{racket} {
	% -- Comments --
	%morecomment=[l]{//},
	%morecomment=[s]{/*}{*/},
	morecomment=[l]{;},         % Inline comments start with ;
	morecomment=[s]{\#|}{|\#},  % Block comments are done with #|  |#
	commentstyle={\color{ForestGreen}\slshape\sffamily},
	% -- Strings --
	morestring=[b]",
	stringstyle=\color{red},
	% --- Literal replacements ---
	literate=%
	%	{lambda}{{$\lambda$}}1
		{->}{{$\rightarrow$}}1
		{*}{{$\times$}}1,
	% -- Keywords (I got these from a vague French website; Racket seems to have no lstlisting style defined anywhere) --
	classoffset=1,
		morekeywords={
			define, define-syntax, define-macro, define-datatype, lambda, define-stream, stream-lambda
		},
		keywordstyle=\color{purple},
	classoffset=2,
		morekeywords={
			begin, call-with-current-continuation, call/cc, call-with-input-file, call-with-output-file, 
			cases, cond, do, else, for-each, if,
			let*, let, let-syntax, letrec, letrec-syntax,
			define-context, define-controller, Integer, Boolean, get, when-required, when-provided,
			maybe_publish, require, submod, or/c, ->, \#\%module-begin, always_publish, with-syntax, define-struct/contract, syntax-case, define/contract,
			let-values, let*-values,
			module, provide,
			and, or, not, 
			delay, force,
			\#`, \#',
			null, cons, car, cdr, cadr, caddr,
			\#lang, implement, begin-for-syntax, rename-out,
			quasiquote, quote, unquote, unquote-splicing,
			map, filter, foldl, syntax, syntax-rules, eval, environment, query 
		},
		keywordstyle=\color{blue},
	classoffset=3,
		morekeywords={import, export},
		keywordstyle=\color{green},
	classoffset=0,
	alsoletter={',`,-,/,>,<,\#,\%},
	moredelim=**[is][\color{lightgray}]{<<@<<}{>>@>>},
	moredelim=**[is][\itshape\color{OliveGreen}]{<<;<<}{>>;>>},
}

% Inline racket
\newcommand{\ir}[1]{\lstinline[basicstyle=\ttfamily\small,language=racket]{#1}}
